% =====================================================================
%  CT Dose AI — Advanced (deep-learning) track: plan and data request
%  Compile with:  pdflatex CT_Dose_AI_Advanced_Track.tex   (run twice)
% =====================================================================
\documentclass[10pt,a4paper]{article}

\usepackage[T1]{fontenc}
\usepackage[utf8]{inputenc}
\usepackage{helvet}\renewcommand{\familydefault}{\sfdefault}
\usepackage{microtype,amsmath,amssymb}
\usepackage[a4paper,top=19mm,bottom=18mm,left=18mm,right=18mm]{geometry}
\usepackage{fancyhdr,titlesec,enumitem,parskip}
\usepackage{graphicx,booktabs,tabularx,array,colortbl}
\usepackage{xcolor,tcolorbox,tikz}
\usepackage[font=small,labelfont=bf,labelsep=period]{caption}
\usepackage[hidelinks]{hyperref}
\usetikzlibrary{arrows.meta,positioning}
\tcbuselibrary{skins,breakable}

\definecolor{ctblue}{HTML}{2A78D6}
\definecolor{ctbluedark}{HTML}{1A4F94}
\definecolor{ctaqua}{HTML}{1BAF7A}
\definecolor{ctorange}{HTML}{EB6834}
\definecolor{ctink}{HTML}{0B0B0B}
\definecolor{ctink2}{HTML}{52514E}
\definecolor{ctline}{HTML}{E2E2DE}
\definecolor{ctsurf}{HTML}{F7F7F5}
\definecolor{ctgreenbg}{HTML}{E8F6F0}

\titleformat{\section}{\Large\bfseries\color{ctbluedark}}{\thesection\ \textperiodcentered\ }{0pt}{}[{\color{ctblue}\titlerule[1.4pt]}]
\titleformat{\subsection}{\large\bfseries\color{ctink}}{\thesubsection}{0.6em}{}
\titlespacing*{\section}{0pt}{15pt}{7pt}
\setlist[itemize]{leftmargin=14pt,itemsep=2pt,topsep=3pt}
\setlist[enumerate]{leftmargin=16pt,itemsep=2pt,topsep=3pt}

\pagestyle{fancy}\fancyhf{}
\renewcommand{\headrulewidth}{0pt}\renewcommand{\footrulewidth}{0.4pt}
\fancyfoot[L]{\footnotesize\color{ctink2}CT Dose AI \textperiodcentered\ advanced track}
\fancyfoot[R]{\footnotesize\color{ctink2}\thepage}

\newcommand{\ctdi}{\ensuremath{\mathrm{CTDI}_{\mathrm{vol}}}}
\newcommand{\mgycm}{mGy$\cdot$cm}

\newtcolorbox{callout}[1][]{enhanced,breakable,colback=ctsurf,colframe=ctblue,boxrule=0pt,leftrule=2.5pt,
  arc=2pt,left=8pt,right=8pt,top=6pt,bottom=6pt,#1}
\newtcolorbox{warnbox}{enhanced,breakable,colback=ctsurf,colframe=ctorange,boxrule=0pt,leftrule=2.5pt,
  arc=2pt,left=8pt,right=8pt,top=6pt,bottom=6pt}
\newtcolorbox{safebox}{enhanced,breakable,colback=ctgreenbg,colframe=ctaqua,boxrule=0pt,leftrule=2.5pt,
  arc=2pt,left=8pt,right=8pt,top=6pt,bottom=6pt}
\newtcolorbox{eqbox}{enhanced,colback=ctsurf,colframe=ctline,boxrule=0.5pt,arc=3pt,
  left=10pt,right=10pt,top=7pt,bottom=7pt}

\newcolumntype{L}{>{\raggedright\arraybackslash}X}
\renewcommand{\arraystretch}{1.25}
\newcommand{\hdr}{\rowcolor{ctsurf}}

\begin{document}
\thispagestyle{empty}

\vspace*{4mm}
{\footnotesize\bfseries\color{ctblue}\MakeUppercase{CT Dose AI \ \textperiodcentered\ \ Advanced track}}

\vspace{3mm}
{\fontsize{24}{28}\selectfont\bfseries Where deep learning earns its place}

\vspace{3mm}
{\large\color{ctink2}Three extensions beyond the physics model: a pre-scan predictor (built), virtual dose
reduction (planned), and a scout-image network (research). With the data each one needs.\par}

\vspace{6mm}
\begin{callout}
\textbf{The premise.} The dose law is physics, so no network will beat it --- that question is settled, and the
report shows the evidence. The open problems are the ones physics cannot answer: \textit{what will the scanner
choose for this patient}, \textit{how low can the dose go before the image stops being diagnostic}, and
\textit{what does this particular body do to the beam}. All three are learnable, and the third one needs images.
\end{callout}

\vspace{2mm}
\noindent\begin{tabularx}{\linewidth}{@{}l L l l@{}}
\toprule
\hdr \textbf{Track} & \textbf{Question it answers} & \textbf{Data needed} & \textbf{Status}\\
\midrule
\textbf{1 \ Pre-scan} & What dose will \textit{this patient} receive, before the scan? & the 50 exams (done) & \textcolor{ctaqua}{\textbf{built}}\\
\textbf{2 \ Virtual dose reduction} & How low can we go, without irradiating anyone to find out? & raw/DICOM images, 30--50 exams & planned\\
\textbf{3 \ Scout network} & What does this body do to the beam? & localizer images, 300+ exams & research\\
\bottomrule
\end{tabularx}

% =====================================================================
\section{Pre-scan dose prediction --- built}

\subsection{The gap it closes}
The dose model needs the mAs. But the tube-current modulation chooses the mAs only when the patient is already on
the table, so before the scan the most useful question --- \textit{what dose will this patient receive?} --- could
not be answered at all. A second, purely statistical stage now answers it:
\begin{eqbox}
\textbf{Stage 1 (learned):} \ BMI $\rightarrow$ mAs the ATCM will choose, \ BMI $\rightarrow$ scan length\\[3pt]
\textbf{Stage 2 (physics):} \ mAs, kV, length $\rightarrow$ \ctdi, DLP
\end{eqbox}

\subsection{Here the flexible models win}
Same competition, same repeated grouped cross-validation as the dose models --- but this time there is no formula
to fall back on, and the ranking inverts:

\begin{table}[h]\footnotesize
\begin{tabularx}{\linewidth}{@{}L r r@{}}
\toprule
\hdr \textbf{Model (BMI as the only input)} & \textbf{mAs error} & \textbf{Scan-length error}\\
\midrule
Median, ignoring BMI (baseline) & 27.9\% & 9.6\%\\
Log-linear in BMI & 23.4\% & 5.6\%\\
Log-quadratic in BMI & 21.4\% & \cellcolor{ctgreenbg}\textbf{5.4\% (selected)}\\
Gaussian process & 21.2\% & 5.3\%\\
$k$-nearest neighbours & 21.4\% & 5.9\%\\
\textbf{Random forest} & \cellcolor{ctgreenbg}\textbf{19.5\% (selected)} & 5.4\%\\
Gradient boosting & 20.6\% & 5.7\%\\
Neural network & 22.8\% & 6.1\%\\
\bottomrule
\end{tabularx}
\end{table}

The random forest beats the linear model by 4 percentage points and the no-BMI baseline by 8. That is a real,
measured gain from a flexible model --- the honest counterpart to the dose result, where the flexible models lost.

\subsection{Honest uncertainty}
Chaining two models means chaining two errors. The app convolves the out-of-fold residual distributions of both
stages (they are independent) rather than adding variances under a normality assumption:

\begin{table}[h]\footnotesize
\begin{tabularx}{\linewidth}{@{}L r r r@{}}
\toprule
\hdr \textbf{Example patient (120 kV)} & \textbf{mAs} & \textbf{\ctdi{} (mGy)} & \textbf{DLP (\mgycm)}\\
\midrule
BMI 22 & 466 \ (309--741) & 28.2 \ (18.7--42.1) & 650 \ (420--1027)\\
BMI 27 & 658 \ (437--1048) & 39.8 \ (26.5--59.5) & 919 \ (593--1452)\\
BMI 35 & 1040 \ (690--1655) & 62.9 \ (41.8--93.9) & 1452 \ (937--2294)\\
\bottomrule
\end{tabularx}
\end{table}

\noindent Roughly $\pm$35\% instead of the $\pm$7\% available once the mAs is known. Wide --- but that width
\textit{is} the information: it says how unpredictable the modulation is, and it is exactly the quantity a
department needs when consenting a patient or forecasting a year's collective dose.

\begin{safebox}
\textbf{Where it is in the app.} A new tab, \textbf{Pre-scan (BMI only)}: expected mAs, \ctdi{} and DLP with
intervals, the two learned curves against the historical exams, and the model competition in an expander. Four new
automated tests cover it, including one that requires the pre-scan interval to be at least twice the width of the
fixed-protocol interval --- so the extra uncertainty can never be quietly lost.
\end{safebox}

% =====================================================================
\section{Virtual dose reduction --- the highest-value next step}

\subsection{The idea}
The clinical bandit learns from radiologist scores, one patient at a time. That is safe but slow: 300 patients to
learn what 300 labels can teach. \textbf{Noise insertion} breaks that limit. Starting from an already-acquired
full-dose scan, one can synthesise the image that \textit{would have been obtained} at 70\%, 50\% or 30\% of the
dose, by adding correlated noise in the projection (or image) domain according to the scanner's noise model.

\begin{callout}
Nobody is irradiated. One existing exam yields four or five labelled examples at different dose levels, and the
radiologist scores images of the \textit{same patient} --- so the comparison is paired, which removes most of the
between-patient variance and multiplies statistical power.
\end{callout}

\subsection{Method}
\begin{enumerate}
\item \textbf{Noise model.} Quantum noise in the projection domain is Poisson in the detected photon count $N$.
To simulate dose factor $f<1$: $\tilde N = \mathrm{Binomial}(N, f)$, or in log domain add Gaussian noise of variance
$\sigma^2 = \frac{1-f}{f}\cdot\frac{1}{N}$, then add the electronic-noise term that dominates at low $N$ and is
what makes very low kV fail. Validate the simulator against two or three real low-dose phantom scans before
trusting it on patients.
\item \textbf{Where.} Projection-domain insertion (needs raw/CT-projection data from the vendor) is the gold
standard; image-domain insertion with a noise-power-spectrum-matched texture is the practical fallback when raw
data is not exportable and is well described in the literature.
\item \textbf{Blinded reading.} Each patient's images at 100/70/50/30\% dose are anonymised, randomly ordered and
scored 1--4 by two readers. Inter-reader agreement ($\kappa$) is reported.
\item \textbf{Fit.} The same Bayesian logistic acceptability model already in the app, now with hundreds of paired
labels; the paired structure is handled with a per-patient random effect.
\item \textbf{Result.} The dose level at which $P(\text{acceptable}) = 0.95$, per BMI band --- measured, on your
own images, with your own readers, before changing a single patient protocol.
\end{enumerate}

\subsection{What to ask for}
\begin{table}[h]\footnotesize
\begin{tabularx}{\linewidth}{@{}l L l@{}}
\toprule
\hdr \textbf{Item} & \textbf{Detail} & \textbf{Who provides}\\
\midrule
Images & 30--50 consecutive CCTA studies, DICOM, anonymised & radiology / PACS\\
Raw data (ideal) & vendor projection data for the same studies, if exportable & vendor / physicist\\
Dose reports & the RDSR of each study (gives true scan length too) & PACS\\
Reading time & 2 readers $\times$ $\sim$4 hours for 150--200 blinded image sets & radiologists\\
Approval & ethics / clinical-audit sign-off for retrospective image reuse & department\\
\bottomrule
\end{tabularx}
\end{table}

\subsection{Effort and output}
Roughly 3--6 weeks of work: simulator (1--2 weeks), validation on phantom (a few days), reading sessions
(scheduled), analysis (1 week). It produces a defensible, site-specific low-dose protocol --- and it is a
publishable study in its own right.

% =====================================================================
\section{Scout-image network --- the research track}

\subsection{Why an image beats BMI}
BMI is a scalar summary of a 3-D body. Two patients with BMI 30 can differ by a factor of two in chest attenuation.
The scanner already acquires the information that matters --- the \textbf{scout / topogram / localizer}, a
low-dose projection radiograph used for planning --- and it is exactly the input the modulation algorithm itself
uses. A convolutional network on the scout can therefore predict what BMI cannot:

\begin{table}[h]\footnotesize
\begin{tabularx}{\linewidth}{@{}l L L@{}}
\toprule
\hdr \textbf{Output} & \textbf{Why it matters} & \textbf{Label source}\\
\midrule
Water-equivalent diameter & enables SSDE --- patient dose, not phantom dose & computed from the axial images (AAPM 220)\\
mAs the ATCM will select & replaces the 19.5\% BMI model; expect single digits & the exam itself\\
Expected image noise & the missing half of the optimisation & measured in a fixed ROI\\
Organ / effective dose & risk, not console metrics & Monte-Carlo simulation as the teacher\\
\bottomrule
\end{tabularx}
\end{table}

\subsection{Architecture and training plan}
\begin{itemize}
\item \textbf{Input:} the AP and lateral scout images, resized to $512\times256$, plus the patient's age, sex and
BMI as tabular side-inputs.
\item \textbf{Backbone:} a small pretrained CNN (ResNet-18 or EfficientNet-B0, ImageNet weights) --- pretraining
matters enormously at these sample sizes --- with the tabular inputs concatenated before the head.
\item \textbf{Heads:} multi-task regression (water-equivalent diameter, mAs, noise), which regularises far better
than three separate models.
\item \textbf{Physics-informed loss:} predict $\log$ mAs and add a consistency term so that the implied
\ctdi{} obeys the fitted scanner law --- the network then cannot contradict the physics it sits on top of. This is
the point at which the word ``physics-informed neural network'' becomes technically correct for this project.
\item \textbf{Validation:} patient-level splits, never image-level; the same repeated cross-validation discipline;
the BMI model as the baseline to beat.
\item \textbf{Data:} 300 exams to show a signal, 1000+ for a solid result. Augmentation (translation, mild scaling,
intensity jitter) and heavy weight decay throughout.
\end{itemize}

\subsection{Deep surrogate for Monte-Carlo dose (stretch goal)}
A full Monte-Carlo organ-dose simulation takes minutes to hours per patient. A network trained on a few thousand
simulated cases reproduces it in milliseconds, which turns organ dose from a research exercise into something the
app can show while the patient is still in the room. This is a genuine use of deep learning as a
\textit{function approximator for an expensive simulator} --- the setting where it is unambiguously the right tool.

\begin{warnbox}
\textbf{Sequencing matters.} Track 3 needs images, an ethics pathway and several hundred exams. Track 2 needs
images from a few dozen exams and produces a clinical result on its own. Track 1 is already running on the data
you have. Do them in that order, and each one funds the next with evidence.
\end{warnbox}

% =====================================================================
\section{The single data request}

Everything above reduces to one email to the department. In priority order:

\begin{enumerate}
\item \textbf{RDSR (radiation dose structured reports)} for the existing 50 exams and, going forward,
automatically. Gives \ctdi, DLP, kV, mAs \textit{and the true scan length} per series --- it settles the
``SCAN SIZE CM'' question and removes all manual typing.
\item \textbf{The protocol fields:} gating mode (prospective/retrospective), collimation, pitch, rotation time,
reconstruction kernel and iterative strength, heart rate.
\item \textbf{A diagnostic-quality score (1--4)} recorded with the report, even informally. This is the label that
makes every learning loop work.
\item \textbf{Anonymised DICOM images} for 30--50 exams, including the scout, for tracks 2 and 3.
\item \textbf{More exams:} 150--300 per exam type is the point where the models stop being limited by sample size.
\end{enumerate}

\begin{safebox}
\textbf{What each request buys.} \ RDSR $\rightarrow$ DLP error from 8\% towards 1--2\%. \
Protocol fields $\rightarrow$ removes the last of the \ctdi{} error and lets the tool reason about gating. \
Quality scores $\rightarrow$ the clinical learning loop starts working. \
Images $\rightarrow$ virtual dose reduction and the scout network. \
More exams $\rightarrow$ the model competition can finally be won by something flexible.
\end{safebox}

\vfill
{\footnotesize\color{ctink2}\rule{\linewidth}{0.4pt}\\[2pt]
CT Dose AI \textperiodcentered\ advanced track \textperiodcentered\ prepared by Thodoris Fragos \textperiodcentered\
repository: \url{https://github.com/Thodoris1657/CT-Dose-AI-Algorithm} \textperiodcentered\ decision-support
software, not a medical device. Any protocol change requires medical-physicist and radiologist approval.\par}

\end{document}
